\section{Discussion and outlook}\label{discussion-and-outlook}

SCORE takes the opposite route to a world model. A world model learns to show only what the player sees, one frame at a time. SCORE builds everything in detail, so that the world cannot break when the player looks somewhere unexpected, and so that it can be used today in existing engines. Neither route is right in general, and world models are improving quickly. What a built world offers is that it can be edited, that it is owned, and that an engine can load it. Its price is that everything has to be built. As models make objects, surfaces and sounds cheap, producing content stops being the hard part, and the creator's direction becomes the scarce input. For that reason SCORE places the creator's decisions at the start of each place, and turns everything that can be checked automatically into checks that need no attention from the creator. Whether this saves the creator time overall has not been established. In one randomised trial with early-2025 tools, experienced developers working on code they knew well were slower with AI assistance than without it \citep{becker2025metr}.

The route has clear weaknesses. Code builders produce clean objects, but an agent has to write a new builder for every kind of object, which scales worse than generation does. A concept picture shows only one side of a place, so the places built from it come out sparser than the concept; the concept-density check catches this but does not fix it, and the optional world step is meant to fix it. Finally, the pace of a build is set by the daily quotas of the picture models, not by the cost of computing.

Beyond the four worlds, SCORE is meant to be packaged so that a creator can run it through a coding assistant. It is also meant to reach further scales: worlds of galaxies or of cells need established open solvers, wrapped as stages of their own, and orbital dynamics through REBOUND \citep{rein2012rebound} is planned as the first of these.
